\documentclass[aspectratio=169]{beamer}
\usetheme{Madrid}
\usecolortheme{seahorse}
\usepackage[utf8]{inputenc}
\usepackage{graphicx}
\usepackage{listings}
\usepackage{hyperref}

\title{ScholarAI: Retrieval-Augmented Generation Learning Assistant}
\subtitle{Node.js + Express + LangChain.js + ChromaDB}
\author{Project Team}
\date{May 17, 2026}

\begin{document}

\begin{frame}
  \titlepage
\end{frame}

\begin{frame}{Problem and Goal}
  \begin{itemize}
    \item Students need fast, reliable answers grounded in course materials.
    \item LLM-only answers can hallucinate or miss citations.
    \item Goal: build a RAG system that returns accurate answers with sources.
  \end{itemize}
\end{frame}

\begin{frame}{What is RAG?}
  \begin{itemize}
    \item Retrieval-Augmented Generation combines search with generation.
    \item Retrieve top-k relevant chunks, then generate with context.
    \item Outcome: higher factuality and traceable citations.
  \end{itemize}
  \vspace{6pt}
  $$\text{Answer} = f(\text{Question}, \text{Retrieved Context})$$
\end{frame}

\begin{frame}{System Architecture}
  \begin{itemize}
    \item Browser UI: upload, chat, sources manager.
    \item Express API: upload, query, sources, health.
    \item RAG engine: load, split, embed, store, retrieve, generate.
  \end{itemize}
  \vspace{4pt}
  \small
  \begin{verbatim}
Browser -> Express Server -> RAG Engine -> ChromaDB + OpenAI
  \end{verbatim}
\end{frame}

\begin{frame}{Core Pipeline: Ingestion}
  \begin{itemize}
    \item Supported files: PDF, DOCX, TXT, MD.
    \item Text extraction: pdf-parse, mammoth, fs.
    \item Chunking: RecursiveCharacterTextSplitter (800 / 150 overlap).
    \item Embeddings: text-embedding-3-small.
    \item Storage: ChromaDB persistent vector store.
  \end{itemize}
\end{frame}

\begin{frame}{Core Pipeline: Query}
  \begin{itemize}
    \item Embed the question using the same embedding model.
    \item Retrieve top-5 chunks by similarity.
    \item Prompt template enforces grounded answers and citations.
    \item LLM: gpt-4o-mini generates final answer.
  \end{itemize}
  \vspace{4pt}
  \small
  \begin{verbatim}
Question -> Embed -> Retrieve -> Prompt -> LLM -> Answer + Sources
  \end{verbatim}
\end{frame}

\begin{frame}{Key Components (LangChain.js)}
  \begin{itemize}
    \item RecursiveCharacterTextSplitter for stable chunking.
    \item OpenAIEmbeddings for vector representation.
    \item Chroma vector store for similarity search.
    \item RetrievalQAChain for end-to-end RAG.
    \item PromptTemplate for structured, cited responses.
  \end{itemize}
\end{frame}

\begin{frame}{API Surface}
  \begin{itemize}
    \item \texttt{GET /api/health}: readiness and source count.
    \item \texttt{GET /api/sources}: list indexed sources.
    \item \texttt{POST /api/upload}: ingest document.
    \item \texttt{POST /api/query}: ask a question.
    \item \texttt{DELETE /api/sources/:name}: remove a source.
  \end{itemize}
\end{frame}

\begin{frame}{Design Decisions}
  \begin{itemize}
    \item Deduplication by SHA-256 hash prevents re-indexing.
    \item Chunk size 800 with 150 overlap balances recall and cost.
    \item text-embedding-3-small for speed and affordability.
    \item RetrievalQAChain returns source docs for citations.
  \end{itemize}
\end{frame}

\begin{frame}{Demo Flow}
  \begin{enumerate}
    \item Upload lecture notes (PDF/DOCX/TXT/MD).
    \item Server extracts and indexes content.
    \item Ask a question in the chat UI.
    \item Receive answer with citations and suggestions.
  \end{enumerate}
\end{frame}

\begin{frame}{Limitations and Future Work}
  \begin{itemize}
    \item Depends on document quality and extraction accuracy.
    \item ChromaDB requires service availability.
    \item Future: per-course namespaces, feedback loops, richer UI analytics.
  \end{itemize}
\end{frame}

\begin{frame}{Thank You}
  \centering
  Questions?
\end{frame}

\end{document}
